% ============================================================================
% CHAPTER 6: HORIZON PLANNING
% ============================================================================

\chapter{Horizon Planning}
\label{ch:planning}

This chapter presents the predictive horizon planner that uses contract composition to anticipate safety violations and enable preemptive controller adaptation.

\section{Predictive Planning Concept}

Traditional reactive control waits for contract violations before adapting. Our horizon-based approach projects current conditions forward in time to predict violations \textit{before they occur}:

\begin{figure}[htbp]
    \centering
    \begin{tikzpicture}[
        scale=0.9,
        transform shape
    ]
        % Timeline
        \draw[-{Stealth}, very thick] (0,0) -- (14,0) node[right] {Time};

        % Current time marker
        \fill[contractpurple] (2,0) circle (0.15);
        \node[below, font=\small] at (2,-0.3) {$t_{\text{now}}$};

        % Prediction horizon
        \fill[contractpurple!30] (2,0.5) rectangle (10,2);
        \node[font=\small] at (6,1.25) {Prediction Horizon};

        % Contract satisfaction curve
        \draw[thick, pidblue] (0,3) -- (2,3) -- (4,2.8) -- (6,2.2) -- (8,1.5);
        \draw[thick, pidblue, dashed] (8,1.5) -- (10,0.8) -- (12,0.5);

        % Threshold line
        \draw[thick, red, dashed] (0,1.5) -- (14,1.5);
        \node[left, font=\small, red] at (0,1.5) {Contract Threshold};

        % Predicted violation
        \fill[red] (8,1.5) circle (0.15);
        \node[above right, font=\small, red] at (8,1.6) {Predicted Violation};

        % Labels
        \node[above, font=\small, pidblue] at (1,3) {Contract Margin};

        % Reactive vs Predictive
        \draw[thick, ->, cbforange] (8,-0.5) -- (8,-1.2);
        \node[below, font=\footnotesize, cbforange] at (8,-1.3) {Reactive: Switch here};

        \draw[thick, ->, hinfgreen] (5,-0.5) -- (5,-1.2);
        \node[below, font=\footnotesize, hinfgreen] at (5,-1.3) {Predictive: Switch early};
    \end{tikzpicture}
    \caption{Predictive vs. reactive controller switching: predictive approach switches before violation}
    \label{fig:predictive_vs_reactive}
\end{figure}

\section{Pacti Library Integration}

The horizon planner leverages the \textbf{Pacti library} for formal contract algebra operations.

\subsection{Pacti Overview}

Pacti is a Python library for assume-guarantee contract reasoning that supports:
\begin{itemize}
    \item Polyhedral constraints (linear inequalities)
    \item Contract composition ($\otimes$)
    \item Contract conjunction ($\wedge$)
    \item Refinement checking ($\preceq$)
    \item Quotient computation ($/$)
\end{itemize}

\subsection{Contract Representation in Pacti}

Contracts are represented as polyhedral sets defined by linear constraints:

\begin{lstlisting}[caption={Pacti contract definition}]
from pacti.contracts import PolyhedralIoContract
from pacti.iocontract import Var

# Define variables
wind_speed = Var('wind_speed')
tracking_error = Var('tracking_error')
controller_ok = Var('controller_ok')

# PID contract in Pacti
pid_contract = PolyhedralIoContract.from_strings(
    input_vars=['wind_speed', 'tracking_error', 'estimation_uncertainty'],
    output_vars=['controller_ok', 'energy_efficiency'],
    assumptions=[
        'wind_speed <= 3.0',
        'tracking_error <= 2.0',
        'estimation_uncertainty <= 1.0'
    ],
    guarantees=[
        'controller_ok >= 1',
        'energy_efficiency >= 0.8'
    ]
)
\end{lstlisting}

\section{Horizon Contract Cascade}

The planner composes contracts over discrete time steps in the prediction horizon:

\begin{figure}[htbp]
    \centering
    \begin{tikzpicture}[
        scale=0.8,
        transform shape,
        contract/.style={
            draw,
            rounded corners=3pt,
            minimum width=2cm,
            minimum height=1.5cm,
            fill=contractpurple!20,
            align=center,
            font=\small
        }
    ]
        % Time steps
        \node[contract] (c0) at (0,0) {$\mathcal{C}(t)$\\Step 0};
        \node[contract] (c1) at (4,0) {$\mathcal{C}(t+\Delta)$\\Step 1};
        \node[contract] (c2) at (8,0) {$\mathcal{C}(t+2\Delta)$\\Step 2};
        \node[contract] (c3) at (12,0) {$\mathcal{C}(t+3\Delta)$\\Step 3};

        % Composition arrows
        \draw[-{Stealth}, thick] (c0) -- (c1) node[midway, above, font=\scriptsize] {$\otimes$};
        \draw[-{Stealth}, thick] (c1) -- (c2) node[midway, above, font=\scriptsize] {$\otimes$};
        \draw[-{Stealth}, thick] (c2) -- (c3) node[midway, above, font=\scriptsize] {$\otimes$};

        % Result
        \node[draw, rounded corners=5pt, fill=hinfgreen!20, minimum width=12cm, minimum height=1cm]
            at (6,-2.5) {Composed Contract: $\mathcal{C}_{\text{horizon}} = \mathcal{C}(t) \otimes \mathcal{C}(t+\Delta) \otimes \mathcal{C}(t+2\Delta) \otimes \mathcal{C}(t+3\Delta)$};
    \end{tikzpicture}
    \caption{Contract cascade over prediction horizon}
    \label{fig:contract_cascade}
\end{figure}

\subsection{Composition Algorithm}

\begin{algorithm}[H]
\caption{Horizon Contract Composition}
\begin{algorithmic}[1]
    \State \textbf{Input:} Current state $\mathbf{x}$, conditions, horizon length $N$, step $\Delta t$
    \State \textbf{Output:} Composed contract, safety margin, recommended controller
    \State
    \State $\mathcal{C}_{\text{composed}} \gets \mathcal{C}_{\text{initial}}$
    \State predictions $\gets []$
    \For{$k = 0$ to $N-1$}
        \State $\mathbf{x}_k \gets$ propagate\_state($\mathbf{x}$, $k \cdot \Delta t$)
        \State conditions$_k \gets$ predict\_conditions($\mathbf{x}_k$, wind\_model)
        \State $\mathcal{C}_k \gets$ build\_step\_contract(conditions$_k$)
        \State $\mathcal{C}_{\text{composed}} \gets \mathcal{C}_{\text{composed}} \otimes \mathcal{C}_k$
        \State predictions.append($\mathcal{C}_k$)
    \EndFor
    \State
    \State safety\_margin $\gets$ evaluate\_margin($\mathcal{C}_{\text{composed}}$)
    \State controller $\gets$ select\_controller(predictions, safety\_margin)
    \State \Return $\mathcal{C}_{\text{composed}}$, safety\_margin, controller
\end{algorithmic}
\end{algorithm}

\section{State Prediction Model}

The planner uses a simplified prediction model for trajectory forecasting:

\subsection{Position Prediction}

\begin{equation}
    \mathbf{p}(t + \tau) = \mathbf{p}(t) + \mathbf{v}(t) \cdot \tau + \frac{1}{2}\mathbf{a}(t) \cdot \tau^2
\end{equation}

\subsection{Wind Effect Prediction}

Assuming constant wind over the horizon:
\begin{equation}
    \mathbf{v}_{\text{eff}}(t + \tau) = \mathbf{v}(t) + \mathbf{v}_{\text{wind}}
\end{equation}

\subsection{Tracking Error Prediction}

\begin{equation}
    e_{\text{track}}(t + \tau) = \|\mathbf{p}_d - \mathbf{p}(t + \tau)\| + \lambda_{\text{dist}} \cdot \max(0, d - d_{\text{threshold}})
\end{equation}

where $\lambda_{\text{dist}}$ penalizes deviation beyond a reasonable distance.

\section{Safety Margin Computation}

The safety margin quantifies how far current conditions are from contract violation:

\begin{equation}
    m_{\text{safety}} = \min_{i} \left( \frac{a_i^{\text{threshold}} - a_i^{\text{actual}}}{a_i^{\text{threshold}}} \right)
\end{equation}

where $a_i$ are the contract assumption variables.

\begin{figure}[htbp]
    \centering
    \begin{tikzpicture}
        \begin{axis}[
            width=10cm, height=6cm,
            xlabel={Time in Horizon (s)},
            ylabel={Safety Margin},
            xmin=0, xmax=3,
            ymin=0, ymax=1,
            grid=major,
            legend pos=south west
        ]
            % Good margin trajectory
            \addplot[thick, hinfgreen, mark=*] coordinates {
                (0, 0.8) (0.5, 0.75) (1.0, 0.7) (1.5, 0.65) (2.0, 0.6) (2.5, 0.55) (3.0, 0.5)
            };

            % Degrading margin trajectory
            \addplot[thick, cbforange, mark=square*] coordinates {
                (0, 0.6) (0.5, 0.5) (1.0, 0.35) (1.5, 0.2) (2.0, 0.1) (2.5, 0.05) (3.0, 0.0)
            };

            % Threshold
            \draw[thick, red, dashed] (axis cs:0,0.2) -- (axis cs:3,0.2);
            \node[font=\footnotesize, red, right] at (axis cs:3,0.2) {Switch Threshold};

            \legend{Nominal, Degraded}
        \end{axis}
    \end{tikzpicture}
    \caption{Safety margin evolution over prediction horizon}
    \label{fig:safety_margin}
\end{figure}

\section{Controller Recommendation Logic}

Based on horizon predictions, the planner recommends a controller:

\begin{lstlisting}[caption={Controller recommendation algorithm}]
def recommend_controller(self, predictions, current_conditions):
    """Recommend controller based on horizon analysis."""

    # Check if any step violates PID assumptions
    pid_viable = True
    for step in predictions:
        if (step['wind_speed'] > 3.0 or
            step['tracking_error'] > 2.0 or
            step['uncertainty'] > 1.0):
            pid_viable = False
            break

    # Check safety-critical conditions
    safety_critical = self._is_safety_critical(predictions)

    if safety_critical:
        return 'Hinf', 0.0  # Force H-inf, zero margin

    if not pid_viable:
        return 'Hinf', self._compute_margin(predictions)

    return 'PID', self._compute_margin(predictions)
\end{lstlisting}

\section{Wind Estimation for Planning}

Accurate wind estimation is critical for meaningful predictions.

\subsection{Low-Pass Filtered Wind Sensor}

\begin{equation}
    \hat{w}(t) = \alpha \cdot w_{\text{sensor}}(t) + (1 - \alpha) \cdot \hat{w}(t - \Delta t)
\end{equation}

where $\alpha = 0.15$ balances responsiveness and stability.

\subsection{Drift-Based Wind Estimation}

When wind sensors are unavailable:
\begin{equation}
    \hat{\mathbf{w}} = \frac{\mathbf{p}_{\text{actual}} - \mathbf{p}_{\text{predicted}}}{\Delta t}
\end{equation}

\importantbox{
The wind filter parameter $\alpha = 0.15$ was tuned empirically. Too high causes oscillation from sensor noise; too low delays response to actual wind changes.
}

\section{Planning Update Cycle}

The horizon planner operates at a slower rate than the control loop:

\begin{table}[htbp]
    \centering
    \caption{Planner timing parameters}
    \label{tab:planner_timing}
    \begin{tabular}{@{}lll@{}}
        \toprule
        \textbf{Parameter} & \textbf{Value} & \textbf{Description} \\
        \midrule
        Update rate & 2 Hz & Planner recalculation frequency \\
        Horizon length & 3.0 s & Lookahead window \\
        Horizon steps & 6 & Discrete steps in horizon \\
        Step size & 0.5 s & Time between horizon steps \\
        \bottomrule
    \end{tabular}
\end{table}

\section{Integrated Planner Architecture}

\begin{lstlisting}[caption={IntegratedHorizonPlanner class structure}]
class IntegratedHorizonPlanner:
    def __init__(self, dt=0.02):
        self.dt = dt
        self.horizon_length = 3.0
        self.horizon_steps = 6
        self.update_interval = 0.5  # 2 Hz

        # Wind estimation
        self.wind_filter_alpha = 0.15
        self.filtered_wind = 0.0

        # Last planning results
        self.last_recommendation = 'PID'
        self.last_safety_margin = 1.0

    def update(self, state, conditions, time):
        """Run horizon planning update."""
        if time - self.last_update_time < self.update_interval:
            return self.last_recommendation, self.last_safety_margin

        # Update wind estimate
        self._update_wind_estimate(conditions)

        # Build horizon predictions
        predictions = self._build_horizon_predictions(state, conditions)

        # Compose contracts
        composed_contract = self._compose_contracts(predictions)

        # Compute safety margin
        safety_margin = self._evaluate_safety_margin(composed_contract)

        # Generate recommendation
        recommendation = self._recommend_controller(predictions, safety_margin)

        self.last_update_time = time
        self.last_recommendation = recommendation
        self.last_safety_margin = safety_margin

        return recommendation, safety_margin
\end{lstlisting}

\section{Planner Contract Specification}

\begin{contractbox}[Horizon Planner Contract $\mathcal{C}_{\text{planner}}$]
\textbf{Assumptions:}
\begin{itemize}
    \item Valid state estimate from EKF
    \item Wind sensor available or drift estimation possible
    \item Sufficient computation time (< 100ms per update)
    \item Previous controller performance data available
\end{itemize}

\textbf{Guarantees:}
\begin{itemize}
    \item Controller recommendation every 0.5s
    \item Safety margin computed for full horizon
    \item Early warning of contract violations (> 1s advance)
    \item Consistent recommendations (no oscillation)
    \item Recommendation valid for current conditions
\end{itemize}
\end{contractbox}

\section{Pre-Flight Mission Verification}

Before takeoff, the planner can verify mission feasibility:

\begin{algorithm}[H]
\caption{Pre-Flight Mission Verification}
\begin{algorithmic}[1]
    \State \textbf{Input:} Waypoints, expected conditions, vehicle capabilities
    \State \textbf{Output:} Mission feasibility, risk assessment
    \State
    \For{each waypoint segment}
        \State Compose segment contracts
        \State Check against vehicle capability contracts
        \State Compute worst-case safety margin
        \If{margin $<$ 0}
            \State \Return INFEASIBLE, segment details
        \EndIf
    \EndFor
    \State
    \State overall\_margin $\gets \min$(segment margins)
    \If{overall\_margin $>$ 0.3}
        \State \Return FEASIBLE, HIGH confidence
    \ElsIf{overall\_margin $>$ 0.1}
        \State \Return FEASIBLE, MODERATE confidence
    \Else
        \State \Return FEASIBLE, LOW confidence (requires H-$\infty$)
    \EndIf
\end{algorithmic}
\end{algorithm}

